\section{Summary of Work Completed}\label{5.1}

This internship set out to build, not merely design, a local-first biology
tutoring system in which a small open-source language model is made
trustworthy by the system around it rather than by its own scale. That
system now exists end to end:

\begin{itemize}
\item A 25.0\,GiB corpus acquired from nine official sources and verified
file-by-file with zero failures.
\item A hybrid dense-and-sparse retrieval system over more than ten
million passages, fused and reranked.
\item A five-way question router covering task type, subdomain, learner
level, teaching style, and retrieval depth.
\item A verification layer combining deterministic recomputation
(arithmetic, genetic-code checks, chemical-equation balancing, pathway
constants) with a learned contradiction detector, calibrated repeatedly
against real generated answers rather than synthetic test cases.
\item An external critic integration that measurably improves on the
deterministic layer alone.
\item A 441-question benchmark, scored by machine-checkable answer keys
rather than by a second model's opinion, used throughout to decide, rather
than assume, whether each addition actually helped.
\end{itemize}

That last discipline produced this internship's most important finding.
The underlying model was not the weak link the project was built to fix:
it already answers the substantial majority of the benchmark correctly
unaided, and it already refuses a fabricated premise correctly seven times
out of eight. A naively built retrieval pipeline can make things
measurably worse, by narrowing what the model is willing to say to only
what a retrieved passage happens to contain. This is not a negative
outcome for the internship. It is the outcome the evaluation methodology
was built to be capable of showing, and it is the substance of the work
reported in Chapter~4.

\section{Current Status}\label{5.2}

The AURA-Bio pipeline is complete and operating end to end: corpus
acquisition and preparation, hybrid dense-and-sparse retrieval, the
five-way question router with its learned classifier, the verification and
critique layers, and the evaluation harness all run against the full
441-question benchmark, with the four-system ablation comparison reported
in full in Section~\ref{4.11}.

The browser interface runs against this pipeline directly. Every page reads
from the local API, and the teaching-quality and failure-review panels
aggregate the pipeline's own traces, so what they display is measured
behaviour rather than illustrative data. A run's routed level and teaching
style, the share of its claims carrying a verifiable citation, its verifier
decision, and any revision it triggered are all recorded per question and
reviewable afterwards.

Per-stage timing of a live request also corrected an assumption worth
recording. Answer generation is the conspicuous cost and was not the
dominant one: verification accounted for 79\% of a request, because the
natural-language-inference checker was being invoked once per
claim-and-passage pair and never using its batch size. Restructuring that
comparison into batched rounds made it roughly seven times faster while
returning byte-identical labels. The component that looks expensive is not
reliably the one that is, and only per-stage measurement distinguishes them.

\section{Fine-Tuning: Plan and Current Progress}\label{5.3}

Fine-tuning the local model is the next major phase of the project. Its
training data has been built and its sizing argued from first principles
rather than guessed at (Section~\ref{5.3.2}), and the adapters are
targeted specifically at the behavioural gap Chapter~4 measures, not at
raw recall, since recall is already close to its ceiling. On that basis,
the grounding adapter is projected to raise the pipeline's claim citation
rate from its currently measured 60.5\% toward the 80--90\% range, an
improvement of roughly +20 to +30 percentage points, with the refusal
examples expected to close most of the remaining gap on adversarial
questions.

\subsection{Why Fine-Tuning Now, and Why This Kind}\label{5.3.1}

Two different training regimes are easy to conflate and require very
different amounts of data and compute: \textbf{continued pretraining}
(also called domain-adaptive pretraining), which trains the model
further on raw domain prose to absorb vocabulary and facts, and
\textbf{instruction tuning}, which trains the model on question-answer
pairs to change its \emph{behaviour}: its format, its citation
discipline, its willingness to refuse a false premise. The measurements in
Chapter~4 point unambiguously at the second. The baseline already scores
highly on raw factual recall. Its most attributable, most measurable gap
is behavioural, and refusal, citation discipline, and level adaptation are
properties learned from demonstrations, not from additional undifferentiated
text.

The economics reinforce the same conclusion. Measured directly on this
corpus, text converts to tokens at approximately 4.1 characters per token,
so one gibibyte of raw prose is roughly 262 million tokens. At a
measured throughput of 200--400 tokens per second for a 4-billion-parameter
model under QLoRA \cite{dettmers2023qlora} on the project's available GTX
1650 hardware, one full epoch over one gibibyte of continued-pretraining
text would take approximately ten to fifteen days of continuous training.
Even using free Kaggle GPU quota, it is roughly thirty to forty-eight hours
of wall-clock time, or about two weeks of quota consumed in short daily
sessions. Spending that budget to move a metric that is already at 0.822
overall, with more than half of all questions already answered perfectly,
is a poor trade against a much smaller, targeted instruction-tuning run
aimed squarely at the measured behavioural gap.

\subsection{Training Data Already Prepared}\label{5.3.2}

An approximately 300\,MB training corpus has been built, mixing
question-answer records with biology prose in a 30:70 ratio, with every
record tagged by provenance so that sources of very different training
value are never silently pooled together.

\begin{table}[h]
\centering
\caption{Instruction-tuning data by provenance}
\label{tab:5.1}
\small
\begin{tabularx}{\textwidth}{@{}lXr@{}}
\toprule
\textbf{Kind} & \textbf{Description} & \textbf{Count} \\
\midrule
\texttt{gold} & Benchmark answers verified against both the answer key and the corpus & 125 \\
\texttt{pubmedqa} & Expert-labelled research questions with evidence and a conclusion & $\sim$723 \\
\texttt{template} & Questions derived mechanically from a chunk's title and section & remainder \\
\bottomrule
\end{tabularx}
\end{table}

Genuine, naturally-occurring question-answer material is scarce (on the
order of 3.5\,MB against the roughly 88\,MB that a 30\% share of a
300\,MB budget requires), so template-derived questions fill the
remainder. These are explicitly \emph{not} to be reported as reasoning
supervision: a question mechanically derived from the answer it is paired
with cannot teach the model to reason toward that answer, only to imitate
its format and citation style. Four content problems in the template
generator were found only by reading actual generated questions rather
than trusting aggregate counts: chemical-identifier lookups masquerading
as concept questions, a severe source imbalance in which one source alone
produced 89\% of all templates before quota-based interleaving was added,
and repeated boilerplate clinical templates that a per-document cap now
suppresses. This is a methodological point (read the content, not the
count) that recurred throughout this project and is reported again here
because it will matter equally during fine-tuning data curation.

Retrieval-optimised sources are deliberately excluded from the training
mix: a UniProt or Gene Ontology record is precise, useful evidence to
retrieve, but training a model on its field-value formatting teaches it to
write like a database record, not like a tutor. Selection instead targets
\emph{teaching register}, meaning textbook-like prose, and a genuine
limitation of the corpus was found in doing so: the available collection
of open textbook material (LibreTexts) totals only 6{,}515 chunks, so even
with maximal weighting toward it, textbook prose reaches only a small
fraction of the training mix at any realistic budget, while Wikipedia and
research abstracts, being effectively unlimited in supply, dominate. This
is recorded as an acquisition gap for future corpus expansion, not a
parameter that further filtering can fix.

\subsection{Planned Adapters and Training Procedure}\label{5.3.3}

Fine-tuning will use QLoRA \cite{dettmers2023qlora}, applying Low-Rank
Adaptation \cite{hu2021lora} to a 4-bit-quantised copy of the base model,
which is the only fine-tuning approach the project's hardware can support
without renting external compute. Three adapters are planned, in order:

\begin{enumerate}
    \item \textbf{Biology teaching adapter}: improves explanation
        structure, worked examples, and level adaptation.
    \item \textbf{Biology grounding adapter}: improves citation
        discipline and, most importantly given the measurements in
        Chapter~4, the model's willingness to refuse a question when the
        available evidence does not support an answer, trained explicitly
        on curated refusal examples rather than left to emerge from prose
        exposure alone.
    \item \textbf{Biology critic adapter} (optional, later): a
        distilled, local replacement for the external critic, once enough
        approved critic judgements have been collected to train one.
\end{enumerate}

The concrete procedure follows six steps, in order:

\begin{enumerate}
    \item Curate several thousand high-quality instruction examples,
        weighted deliberately toward refusal and grounding demonstrations
        rather than pure factual recall.
    \item Hold out a strict evaluation split that never touches training.
    \item Configure and run a QLoRA pilot at a small scale first.
    \item Compare the adapter against the base model on the existing
        441-question benchmark under the \texttt{finetuned\_rag} system
        slot already built into the evaluation harness (Section~\ref{3.8}).
    \item Explicitly check for catastrophic degradation on questions the
        base model already answered correctly, since a fine-tune that
        fixes hallucination while breaking factual recall would not be a
        net improvement.
    \item Only then integrate adapter selection into the routing layer.
\end{enumerate}

An initial instruction-tuning target of five thousand to twenty thousand
curated examples is planned, consistent with the sizing argument above,
with quality prioritised over reaching any particular quantity.

\section{Future Scope: Generalising Beyond Biology}\label{5.4}

Biology was chosen as the first domain because it is broad enough to
exercise every stage of the architecture while remaining a single,
well-scoped subject for an initial build. Nothing in the architecture
itself (the router, the hybrid retrieval stack, the verification
composition pattern, the critic interface, or the benchmark methodology)
is biology-specific in its \emph{shape}. What is biology-specific is the
content plugged into each stage: the corpus, the router's subdomain
labels, the specific deterministic checkers, and the benchmark questions.
The long-term vision for this project is therefore not to build a second,
separate system for each new subject, but to factor out a \textbf{shared
core} and let each new subject be added as a self-contained plugin.

\subsection{A Modular Plugin Interface}\label{5.4.1}

Each future domain module is intended to define the same seven artefacts
biology already has, and nothing else:

\begin{center}
\footnotesize\ttfamily
\begin{tabular}{@{}ll@{}}
domain name & e.g.\ ``coding'', ``mathematics'', ``physics'' \\
corpus & domain-specific retrieval material, acquired and verified \\
router labels & task types and subdomains specific to the field \\
prompt templates & teaching-mode structures for the subject \\
verification tools & the deterministic checkers that make sense here \\
evaluation dataset & a benchmark with machine-checkable answer keys \\
optional LoRA adapter & a domain-adapted version of the base model \\
dashboard metrics & what the ablation and evaluation pages show \\
\end{tabular}
\end{center}

The core insight carried over from biology is that \textbf{a checker
should recompute, not merely pattern-match}, wherever the domain allows
it, exactly as SymPy recomputes a genetics ratio instead of reading it
off the model's stated working, and exactly as Biopython checks a codon
against the actual genetic code table instead of trusting the model's
claim about it. This principle turns out to generalise unusually well to
other STEM subjects, several of which are naturally \emph{more}
amenable to deterministic checking than biology is.

\subsection{Candidate Modules}\label{5.4.2}

\textbf{Mathematics.} The most natural next module, since SymPy
\cite{meurer2017sympy} is already a core dependency of the biology
verifier. A mathematics module would parse a stated equation or proof step
and re-derive it symbolically, checking not just a final numeric answer
but the validity of each intermediate step, a strictly stronger form of
verification than anything currently possible for a free-text biology
explanation, because mathematical claims are, unlike most biological ones,
usually fully formalisable.

\textbf{Physics.} A physics module would carry a formula database, a
dimensional-analysis checker (verifying that a stated quantity's units are
even physically possible for what it claims to measure, directly
analogous to the unit-magnitude checker already built for biology,
Section~\ref{3.6}), and a numerical solver for standard problem types.

\textbf{Coding.} A coding module changes what ``retrieval'' and
``verification'' mean rather than removing them: the corpus becomes source
repositories and documentation, chunking becomes code-aware (splitting on
function and class boundaries rather than token windows), and the
verification layer becomes an actual runtime sandbox running the
project's existing test suite or a generated one via \texttt{pytest},
plus static linting. This is verification that does not approximate
correctness but \emph{executes} it, which is an even stronger guarantee
than anything deterministic recomputation offers in biology.

\textbf{Logic.} A logic module would use constraint solvers and symbolic
reasoning engines to check multi-step deductive arguments for internal
consistency, and to detect contradictions across a chain of stated
premises, a symbolic, guaranteed-sound counterpart to the statistical
natural-language-inference contradiction check built for biology
(Section~\ref{4.6.3}).

\subsection{Why This Specific Order}\label{5.4.3}

Mathematics and physics are prioritised over coding and logic in the
near term because they reuse SymPy and the dimensional-analysis pattern
that already exist, and because their claims are more often fully
verifiable by recomputation, the strongest and cheapest form of
verification this project has found and the one every measurement in
Chapter~4 shows is worth building wherever it is possible. Coding is the
module where the architecture generalises furthest beyond biology's
pattern, since ``does this code actually run and pass its tests'' is a
categorically stronger check than anything text-based verification can
offer, and is a natural medium-term target once the shared core is
factored out cleanly enough to support it.

\section{Final Positioning}\label{5.5}

AURA-Bio should be understood as a transparent, evidence-grounded,
adaptive biology tutoring system that combines a lightweight local
language model with a curated corpus, hybrid retrieval, learner-aware
pedagogy routing, biological verification, structured critique, and
evaluation-driven improvement. It should not be understood as a biology
chatbot, as a general-purpose conversational assistant with a biology
skin, or as a fully autonomous teaching system that replaces a human
instructor. Its central contribution, argued and measured throughout this
report, is that a lightweight local model can be made more reliable, more
adaptive, and more educationally useful through system-level architecture
rather than through model scale alone, and that this claim is only worth
making when it is backed by a benchmark honest enough to show where the
architecture is helping and where it is not.

\section{Concluding Remarks}\label{5.6}

The most valuable outcome of this internship was not a single accuracy
number. It was building an evaluation methodology rigorous enough to show,
repeatedly and specifically, when an intuitively reasonable design
decision was quietly making the system worse, and then having the
discipline to trust that measurement over the intuition. A small language
model, it turns out, already knows a great deal of biology. What it cannot
do on its own is show its working, admit what it does not know, or prove
that it checked its own arithmetic. Building the system that lets it do
all three, and proving with real numbers that each piece of that system
earns its place, is the work this report describes, complete end to end
and clearly scoped, in Sections~\ref{5.3} and~\ref{5.4}, for what comes
next.
